\section*{Part X: The Founder's Field Manual}
\addcontentsline{toc}{section}{Part X: The Founder's Field Manual}

\subsection{Worksheet A: Process Design Dossier}
\label{sec:worksheet_a_process_dossier}

For 1 critical workflow, record:
\begin{enumerate}
    \item \textbf{Customer outcome:} the tangible value delivered to the recipient;
    \item \textbf{Unit of work:} bounded transactional entity with clear start and end states;
    \item \textbf{Start event and trigger:} what officially initiates active work;
    \item \textbf{Acceptable completion definition:} verifiable conditions for successful closure;
    \item \textbf{Customer touchpoints:} every point where the user interacts or waits;
    \item \textbf{Internal activities and owners:} sequential tasks and named role accountabilities;
    \item \textbf{Inputs and dependencies:} prerequisites, data assets, and required tooling;
    \item \textbf{Quality checkpoints:} gates where defects are blocked from proceeding;
    \item \textbf{Normal cycle time and rework patterns:} expected duration and known failure modes;
    \item \textbf{Escalation criteria and record location:} authoritative log and trigger thresholds.
\end{enumerate}

Add the date of last actual observation rather than merely a process-map approval date. Use a real case to test each field. If ``the process owner'' cannot identify the oldest incomplete item, the design is not operating as written. Begin with 1 process that materially affects customer trust or founder workload. Do not attempt to document the entire organization in 1 afternoon.

\subsection{Worksheet B: Quality and Risk Register}
\label{sec:worksheet_b_risk_register}

Create columns for failure mode, consequences, customer visibility, likelihood estimate, detectability, prevention control, detection control, recovery path, and owner. Use qualitative ratings when evidence is scarce; do not multiply invented precision scores and treat the resulting risk rank as scientific. Identify high-severity failures that require mandatory action irrespective of calculated frequency.

The worksheet should surface cases like inaccurate completion claims, exposed learner records, late contractual delivery, or unverified AI feedback. Revisit the register after real incidents. A risk register whose last update predates the most recent operating change is not a functioning control.

\subsection{Worksheet C: Operational Metrics Dictionary}
\label{sec:worksheet_c_metrics_dictionary}

For every metric, specify:
\begin{itemize}
    \item Core purpose and decision it informs;
    \item Explicit formula and mathematical definition;
    \item Event start and stop boundaries;
    \item Denominator definition and active exclusions;
    \item Primary data source and automated or manual collection method;
    \item Accountable metric owner and review cadence;
    \item Target rationale and historical baseline;
    \item Anticipated gaming behavior and counter-measures;
    \item Balancing paired measure (such as speed paired with audit accuracy).
\end{itemize}

Record baseline and intervention dates so improvements can be compared fairly. If a metric changes definition, keep the historical version or explicitly restate prior values. Avoid using individual productivity comparisons without accounting for case difficulty, training responsibilities, and unusual conditions.

\subsection{Worksheet D: Bottleneck and Capacity Calculator}
\label{sec:worksheet_d_capacity_calculator}

List each work stage, volume per period, effective service minutes per unit, resource hours available, utilization, queue, and quality constraint. Effective hours should reflect meetings, breaks, training, coordination, and ordinary interruptions. Then calculate where demand exceeds capacity and what variation could do to waiting time.

A useful scenario table includes low, expected, and peak demand; ordinary and high-complexity cases; reviewer absence; and rework. State which assumptions are measured and which are guesses. Capacity numbers without quality and timing commitments are not meaningful. Do not silently count the founder's unsustainable overtime as free spare capacity.

\subsection{Worksheet E: Automation Investment Memo}
\label{sec:worksheet_e_automation_memo}

The 1-page memo should state:
\begin{enumerate}
    \item Current manual workflow and baseline labor minutes per transaction;
    \item Historical defect rate and downstream rework costs;
    \item Proposed automation boundary and operational posture (assistance vs. autonomy);
    \item Required input datasets, data classifications, and privacy controls;
    \item Identified failure modes, adversarial risks, and silent degradation risks;
    \item Intended human oversight mechanism, reviewer workload, and rejection path;
    \item Implementation, prompt engineering, integration, and ongoing maintenance cost;
    \item Expected net time and quality impact evaluated under full system cost;
    \item Independent validation test harness and passing thresholds;
    \item Rollback procedure, manual fallback, and stop-the-line criteria.
\end{enumerate}

Approve a bounded pilot, not permanent autonomy on the strength of a demonstration. Evaluate realistic representative cases, including rare but costly errors. A tool that makes low-value output faster but increases scarce reviewer workload is not an operational gain.

\subsection{Worksheet F: Knowledge-Asset Specification}
\label{sec:worksheet_f_knowledge_asset}

For every reusable asset, record why it exists, who uses it, owner, source and rights, authoritative location, version, review date, contexts where it applies, contexts where it does not, related examples, and retirement criteria. Label confidential and personal material appropriately. A useful knowledge asset should be usable by somebody who did not create it.

For an educational rubric, attach anonymized exemplar work and rater calibration notes. For an AI evaluation dataset, document provenance, limitations, and whether it may contain sensitive or restricted information. For a college contract template, identify what must be locally reviewed rather than copied uncritically.

\subsection{Worksheet G: Weekly Operations Review}
\label{sec:worksheet_g_weekly_review}

Use 1 page with: commitments due, aging work, customer-impacting incidents, quality exceptions, verified outcomes, staff capacity, cash-relevant bottlenecks, and improvement experiments. Review the few material deviations. Assign actions with owners and dates. Next week, begin by checking whether the prior decisions actually changed the system.

The ideal review is short because the underlying information is reliable. If it repeatedly becomes a debate over basic facts, fix data capture and ownership before adding more KPIs. If a high-quality weekly review yields no decisions for several weeks, consider whether the cadence is necessary.

\subsection{Worksheet H: A 30-Day Operational Excellence Sprint}
\label{sec:worksheet_h_30_day_sprint}

A structured 4-week sprint transforms a bottleneck into standard work:
\begin{itemize}
    \item \textbf{Week 1 (Diagnosis):} choose a customer-critical workflow, map 5 actual cases, and measure touch time, waiting, rework, and exceptions;
    \item \textbf{Week 2 (Standards and Gates):} set an acceptance standard, define 1 or 2 meaningful measures, and repair the largest known handoff failure;
    \item \textbf{Week 3 (PDSA Pilot):} run a bounded PDSA experiment and inspect quality, waiting, and human burden under realistic arrival;
    \item \textbf{Week 4 (Codification and Handover):} document the resulting standard, assign ownership, and write a short decision memo about what remains uncertain.
\end{itemize}

The outcome should be 1 functioning, improved process, not a large binder. If no paid customers exist yet, use a realistic rehearsal, mark its limitations, and avoid claiming that simulated performance proves external service reliability.

\subsection{Worksheet I: Operational Capability Audit}
\label{sec:worksheet_i_capability_audit}

Score each critical workflow descriptively on 5 dimensions:
\begin{enumerate}
    \item \textbf{Outcome clarity:} observable, verified definition of finished work;
    \item \textbf{Repeatability:} consistent execution across different contributors;
    \item \textbf{Observability:} visible queues, error rates, and bottlenecks;
    \item \textbf{Recovery:} disciplined response, containment, and learning from incidents;
    \item \textbf{Transferability:} ability of a new competent colleague to operate the workflow without founder assistance.
\end{enumerate}

The purpose is to locate a constraint, not to declare a grand organization-wide maturity level. Require evidence: completed cases, dated logs, actual incident resolutions, 2nd-person execution, and measured results. An operation can be strong in repeatability yet weak in adaptability, or strong in customer intimacy yet dependent on the founder. Define the next improvement in observable terms: \textit{``another reviewer can use the rubric to reach a defensible decision without asking the founder''} is a valid test.

\subsection{12 Recurring Errors in Building Operations}
\label{sec:twelve_recurring_errors}

Common operational pitfalls include:
\begin{enumerate}
    \item Optimizing an output that customers do not value;
    \item Treating average workload as peak capacity;
    \item Confusing fast ticket closure with actual problem resolution;
    \item Standardizing legitimate customer differences out of existence;
    \item Documenting elaborate procedures that nobody actually follows;
    \item Automating a workflow before defining its operational boundaries;
    \item Counting unverified AI drafts as completed, validated work;
    \item Hiding founder overtime to make unit economics look artificially attractive;
    \item Rewarding a single metric while ignoring paired quality safeguards;
    \item Blaming individual workers for predictable system and handoff failures;
    \item Letting operational incidents disappear without verified corrective actions;
    \item Accepting commercial contracts that the delivery system cannot reliably honor.
\end{enumerate}

Each failure has a similar underlying mechanism: a representation of the work becomes detached from the actual customer consequence. The remedy is direct observation, explicit definitions, proportionate ownership, and the courage to revise the system when evidence contradicts the founder's preferred story.
